\section{Strategies excluded, with witnesses}\label{app:graveyard}

By Corollary~\ref{cor:nomargin} the selection inequality
\eqref{eq:dominate} is an equality on a set of full dimension in the
weights.  A device whose output is a strictly positive quantity therefore
asserts something false at every point of that set.  Four of the seven
strategies below fail for that reason alone.

For a $k$-subset $C$ of a design write $\Gram_C$ for the Gram matrix
$(\langle g_c,g_d\rangle)_{c,d\in C}$.  By \eqref{eq:twoproducts},
$\Gram_C=M_C^{\vphantom{\mathsf{H}}}M_C\adj$ and
$\SC=M_C\adj M_C^{\vphantom{\mathsf{H}}}$ with
$M_C$ square, so the two have the same spectrum and
\begin{equation}\label{eq:tieleg}
  \delta(C)\;:=\;\det\bigl(\Gram_C-I_k\bigr)\;=\;\det\bigl(\SC-I_k\bigr) .
\end{equation}
Domination of $C$ implies $\delta(C)\ge0$; we call $\delta$ the
\emph{determinant leg}.  At rank three we also use
\begin{equation}\label{eq:twomoment}
  \gamma(C)\;:=\;4\det\Gram_C-\bigl(\tr\Gram_C-1\bigr)^2 ,
\end{equation}
a function of the first two moments of the spectrum of $\SC$ alone.

\subsection{Certificates with a uniform floor}\label{app:certificates}

The natural request is a nonnegative combination of the determinant legs,
at each design, bounded below by a positive constant.  No such
combination exists, and the obstruction has nothing to do with degree,
polynomiality or equivariance: the coefficients below may be arbitrary
nonnegative functions of the design.

\begin{proposition}[Lean]\label{prop:noaggregate}
There is no $\theta>0$ with the following property: for every design $\dsn$ at
$(7,3)$ with $\ell_c>1$ for all $c$ there are numbers $w_C(\dsn)\ge0$, indexed by
the $35$ triples, with
\[
  \sum_{\lvert C\rvert=3} w_C(\dsn)\,\delta(C)\;\ge\;\theta .
\]
\end{proposition}

\begin{proof}
Let $\dsn_{\triangle}$ be the split tetrahedron at $(7,3)$ of \S\ref{app:split}: seven
atoms carrying the four tetrahedron directions with multiplicities $2,2,2,1$
and weights $(\tfrac18,\tfrac18,\tfrac18,\tfrac14,\tfrac18,\tfrac18,\tfrac18)$.
Every leverage equals $3$, so $\dsn_{\triangle}$ lies in the stated class.  A triple of
$\dsn_{\triangle}$ carries either three distinct directions or two equal ones.  In the
first case $\Gram_C$ has diagonal $3$ and off-diagonal $-1$, so
$\Gram_C-I_3=3I_3-J_3$, of eigenvalues $0,3,3$ and determinant $0$.  In the
second, after ordering so that the first two atoms agree,
\[
  \Gram_C-I_3=\begin{pmatrix}2&3&-1\\ 3&2&-1\\ -1&-1&2\end{pmatrix},
  \qquad \det=2(4-1)-3(6-1)+(-1)(-3+2)=-8 .
\]
Hence $\delta(C)\le0$ for all $35$ triples, and every nonnegative combination
of the legs at $\dsn_{\triangle}$ is at most $0$.
\end{proof}

Any Positivstellensatz certificate in Putinar or Schmüdgen form yields, on
evaluation at a design, exactly such a combination, its coefficients being the
sums of squares evaluated there and its floor the constant of the
representation.  So:

\begin{corollary}\label{cor:noputinar}
No representation
$\sum_{C}\sigma_C\,\delta(C)=1+\sigma_0$, valid on the all-heavy locus at
$(7,3)$ with $\sigma_0$ and the $\sigma_C$ nonnegative there, exists --- at any
degree.
\end{corollary}

\begin{proof}
Evaluate at $\dsn_{\triangle}$: the left side is at most $0$ and the right side at least
$1$.
\end{proof}

Consider the two systems on the all-heavy locus,
\[
  \text{(strict)}\quad \delta(C)<0 \ \ \text{for all } C,
  \qquad\qquad
  \text{(closed)}\quad \delta(C)\le0 \ \ \text{for all } C .
\]
Failure of $\Gtz(7,3)$ would give a point of the strict system; the design
$\dsn_{\triangle}$ is a point of the closed one.  So the closed system is
inhabited and carries no certificate of emptiness of any kind, while the
strict system is the one that must be certified empty.  The two systems
differ on the equality locus, and nowhere else.

\begin{remark}\label{rem:stengle}
Stengle's form, in which the multiplicative monoid part is nonempty,
survives; one can bound its size.  A Stengle certificate of emptiness of
the strict system is an identity $f+g^2+h=0$ with $f$ in the cone generated by
the defining inequalities, $h$ in the ideal generated by the equations, and
$g$ a product of powers of the strict generators $-\delta(C)$.  At any point of
the closed system all generators are nonnegative, so $f\ge0$ and $h=0$ there,
forcing $g=0$.  Now the closed system contains one split tetrahedron for each
partition of the seven indices into classes of sizes $2,2,2,1$, of which there
are $7!/(2!^3\,3!)=105$; at the partition $\Pi$ the leg $\delta$ vanishes
exactly on the $20$ triples meeting three distinct classes of $\Pi$.  So the set
of triples occurring as factors of $g$ must meet, for each of the $105$
partitions, that partition's own family of $20$.  No set of at most three
triples meets all $105$ families; $5355$ sets of four do, among them
$\{123,124,135,145\}$.  Since $\delta$
has degree $3$ in the entries of $\Gamma$ and degree $6$ in the atom
coordinates, the monoid part of every Stengle certificate has degree at least
$12$ in the $28$ Gram entries and at least $24$ in the $21$ coordinates.  A sum
of squares of degree $12$ in $28$ variables has a half-degree monomial basis of
$\binom{34}{6}=1\,344\,904$ elements.
\end{remark}

\subsection{Margin transfer across a deletion}\label{app:transfer}

Induction on the size deletes an atom, applies $\Gtz(m-1,k)$, and lifts.
The lift is exact, and one computation settles what it costs.

Let $\dsn$ be a design at $(m,k)$ and $e$ an index with $s_e<1$, equivalently
$I-t_e\,g_e^{\vphantom{\mathsf{T}}}g_e\tp\pd0$.  Choose $R$ with
$R\tp\bigl(I-t_e g_e^{\vphantom{\mathsf{T}}}g_e\tp\bigr)R=I$ and set
\[
  g'_c:=\sqrt{1-t_e}\;R\tp g_c,
  \qquad
  t'_c:=\frac{t_c}{1-t_e}
  \qquad (c\ne e).
\]
The weights are positive and sum to one, and
\[
  \sum_{c\ne e}t'_c\,g'^{\vphantom{\mathsf{T}}}_c{g'_c}\tp
  =R\tp\Bigl(\sum_{c\ne e}t_c\,g_c^{\vphantom{\mathsf{T}}}g_c\tp\Bigr)R
  =R\tp\bigl(I-t_e g_e^{\vphantom{\mathsf{T}}}g_e\tp\bigr)R=I,
\]
so $\dsn'$ is a design at $(m-1,k)$.

\begin{proposition}\label{prop:deletionprice}
Let $C\subseteq\{c:c\ne e\}$ have $\lvert C\rvert=k$ and satisfy
$S'_C\psd(1+\varepsilon)I$ in $\dsn'$, with $\varepsilon\ge0$.  Then $C$
dominates in $\dsn$ provided
\begin{equation}\label{eq:kappa}
  1+\varepsilon\;\ge\;\kappa_e^{-1},
  \qquad\text{where}\qquad
  \kappa_e:=\frac{1-s_e}{1-t_e} ,
\end{equation}
that is provided $\varepsilon\ge t_e(\ell_e-1)/(1-s_e)$.  The threshold is at
most $0$ if and only if $\ell_e\le1$.
\end{proposition}

\begin{proof}
From $R\tp\bigl(I-t_eg_e^{\vphantom{\mathsf{T}}}g_e\tp\bigr)R=I$ one reads
$RR\tp=\bigl(I-t_eg_e^{\vphantom{\mathsf{T}}}g_e\tp\bigr)^{-1}$, hence for
$y:=R^{-1}w$,
\[
  y\tp S'_Cy=(1-t_e)\,w\tp S_Cw,
  \qquad
  y\tp y=w\tp\bigl(I-t_eg_e^{\vphantom{\mathsf{T}}}g_e\tp\bigr)w
        =\lVert w\rVert^2-t_e\langle g_e,w\rangle^2 ,
\]
the first because $S'_C=(1-t_e)R\tp S_CR$.  Applying the hypothesis at $y$ and
taking $\lVert w\rVert=1$,
\[
  w\tp S_Cw\;\ge\;\frac{(1+\varepsilon)\bigl(1-t_e\langle g_e,w\rangle^2\bigr)}
                       {1-t_e}
   \;\ge\;\frac{(1+\varepsilon)(1-s_e)}{1-t_e}
   \;=\;(1+\varepsilon)\,\kappa_e ,
\]
by Cauchy--Schwarz, which gives $\langle g_e,w\rangle^2\le\ell_e$ and hence
$t_e\langle g_e,w\rangle^2\le t_e\ell_e=s_e$.  This is at least $1$
exactly under \eqref{eq:kappa}.  Finally
$\kappa_e^{-1}-1=(s_e-t_e)/(1-s_e)=t_e(\ell_e-1)/(1-s_e)$, whose sign is that
of $\ell_e-1$.
\end{proof}

Proposition~\ref{prop:deletionprice} is an identity, and it splits the
induction into the two branches already in use.  At $\ell_e\le1$
the threshold is nonpositive and the deletion needs nothing downstairs;
this is the light-atom step that reduces $\Gtz(m,k)$ to all-heavy
designs.  At $\ell_e>1$ the threshold is strictly positive.

\begin{corollary}\label{cor:noheavydeletion}
At rank three no deletion of a heavy atom transports domination without a
margin, and no such margin is available: by Corollary~\ref{cor:nomargin}, for
every $m\ge5$ and every $\varepsilon>0$ there is a design at $(m-1,3)$ with no
$k$-subset satisfying $\SC\psd(1+\varepsilon)I$.
\end{corollary}

No choice of the deleted atom repairs this, because the threshold depends on
$e$ only through $\ell_e$ and by Lemma~\ref{lem:trace} some atom has
$\ell_e\ge k$.

\subsection{Restricting the region}\label{app:region}

Cutting the degenerate part out of the design space and running a
compactness argument on the rest requires a cut that misses the equality
locus.  The two cuts that suggest themselves, a floor on the weights and
the exclusion of parallel atoms, both leave the locus intact.

By Theorem~\ref{thm:tieseverywhere} every weight vector on $m\ge k+1$ indices
carries an extremal design, one for each partition of the indices into $k+1$
nonempty classes; the construction is displayed in \S\ref{app:split}.  So no
restriction of the weight simplex --- to a compact subset, to a
neighbourhood of the uniform point, to any nonempty relatively open set ---
excludes the equality locus, and the infimum of the selection objective over
any such region is at most $1$.

The second cut fails for a different reason.  Every design produced by
splitting has two parallel atoms, so excluding parallel pairs removes the
constructed part of the equality locus.  It does not remove all of it: the
diamond of Example~\ref{ex:diamond}, whose data is in \S\ref{app:diamond}, is
an extremal design at $(5,3)$ whose five atoms are pairwise non-parallel.  Its
splittings give extremal designs with a parallel pair at $(6,3)$ and $(7,3)$,
and perturbing a split pair apart moves the design off the parallel locus while
moving the value by an arbitrarily small amount, since the value is continuous.

\subsection{The compression lift at positive weight}\label{app:lift}

The lift \eqref{eq:schur} used in Theorem~\ref{thm:boundary} holds at $t_z=0$
by an exact identity and fails at every positive $t_z$, however small
(Theorem~\ref{thm:sharp}).  The witness is the chart point and one-parameter
weight family of \S\ref{app:liftfail}, whose relevant $2\times2$ block has
determinant $-\tau(2+3\tau)/9$ at pivot weight $\tau$.  The failure is not a
matter of size: the deficit is of first order in $\tau$ against a downstairs
slack which the family freezes at $0$.  This is again
Corollary~\ref{cor:nomargin}, so the boundary analysis of
\S\ref{sec:boundary} runs at the limit point.

\subsection{Real-rootedness of the mixture}\label{app:rooted}

Theorem~\ref{thm:average} makes the selection problem a derandomization:
$\E_\pi[\SC]\psd I$, and one wants a single $C$.  The interlacing method
supplies derandomizations of just this shape, and the mixture
$q=\E_\pi[\chi_{\SC}]$ of Proposition~\ref{prop:qone} is the object it
would act on.  Real-rootedness of $q$ alone, however, buys nothing.

\begin{proposition}\label{prop:rootedness}
There are real-rooted monic quadratics $f_1,f_2$ whose average is real-rooted
and whose average has least root strictly larger than the least roots of both.
\end{proposition}

\begin{proof}
Take
\[
  f_1=X^2-X,\qquad f_2=X^2-23X+60,
\]
with roots $\{0,1\}$ and $\{3,20\}$, so least roots $0$ and $3$.  Their average
is $X^2-12X+30$, of discriminant $144-120=24>0$, with roots $6\pm\sqrt6$; the
least is $6-\sqrt6=3.5505\ldots>3$.
\end{proof}

The hypothesis the interlacing method actually requires is that the whole
convex family be real-rooted --- a common interlacing --- and the pair above
fails it:
\[
  \lambda f_1+(1-\lambda)f_2=X^2+(22\lambda-23)X+60(1-\lambda),
  \qquad
  \operatorname{disc}=484\lambda^2-772\lambda+289 ,
\]
which at $\lambda=\tfrac45$ equals $-\tfrac{471}{25}<0$.  So the conclusion of
\cite{MSS2015} is unavailable for this pair, and no contradiction
arises.  Nor is real-rootedness
preserved by averaging in general: $X^2-1$ and $X^2+10X+25$ are real-rooted and
their average $X^2+5X+12$ has discriminant $-23$.  A derandomization of
Theorem~\ref{thm:average} must therefore verify the interlacing hypothesis for
volume sampling, not assume real-rootedness of the mixture.  This is
Problem~\ref{prob:interlace}.

\subsection{Two moments}\label{app:moments}

At rank three there is a genuine criterion in the first two moments of the
spectrum.  It is sharp among such criteria, and it does not decide the open
cells.

\begin{proposition}[Lean]\label{prop:twomoment}
Let $C$ be a $3$-subset of a design at rank three with
$\ell_C:=\tr\Gram_C>3$.  If $\gamma(C)\ge0$ then $C$ dominates.
\end{proposition}

\begin{proof}
Let $x\le y\le z$ be the eigenvalues of $\SC$, all nonnegative, with
$x+y+z=\ell_C$ and $xyz=\det\Gram_C$.  By the arithmetic--geometric mean
inequality applied to $y,z$,
\[
  4\det\Gram_C=4xyz\le 4x\Bigl(\frac{y+z}{2}\Bigr)^2=x(\ell_C-x)^2 .
\]
The hypothesis $\gamma(C)\ge0$ reads $4\det\Gram_C\ge(\ell_C-1)^2$, so
$x(\ell_C-x)^2-(\ell_C-1)^2\ge0$.  Expanding both sides,
\[
  x(\ell_C-x)^2-(\ell_C-1)^2=(x-1)\bigl(x^2+(1-2\ell_C)x+(\ell_C-1)^2\bigr),
\]
an identity in $x$ and $\ell_C$.  Write $Q(x)$ for the quadratic factor.  Its vertex
is at $x=\ell_C-\tfrac12>1$, so $Q$ decreases on $[0,1]$, and
$Q(1)=2-2\ell_C+(\ell_C-1)^2=(\ell_C-1)(\ell_C-3)>0$; hence $Q>0$ on $[0,1]$.  If $x<1$ then
$(x-1)Q(x)<0$, contradicting the displayed inequality.  So $x\ge1$, that is
$\SC\psd I$.
\end{proof}

The quadratic factor $Q$ is the multiplier of the certificate; no
semidefinite programming enters anywhere.  At the tetrahedron
(\S\ref{app:tetra}) one has $\ell_C=9$, $\det\Gram_C=16$ and
$\gamma(C)=0$, so the criterion is tight on the locus it was built
to reach: the arithmetic--geometric step is an equality exactly when the
two upper eigenvalues coincide, as they do at $\{1,4,4\}$.

\begin{proposition}\label{prop:momentsharp}
Let $\ell_C>3$ and $\varepsilon\ge0$ satisfy $4\varepsilon<(\ell_C-1)^2$.
There is a positive semidefinite $3\times3$ matrix of trace $\ell_C$ and
determinant $\varepsilon$ whose least eigenvalue is smaller than $1$.  Hence
no criterion depending on $(\tr\Gram_C,\det\Gram_C)$ alone is stronger than
Proposition~\ref{prop:twomoment}.
\end{proposition}

\begin{proof}
On $[0,1]$ the function $x\mapsto x(\ell_C-x)^2$ has derivative $(\ell_C-x)(\ell_C-3x)>0$, so
it increases from $0$ to $(\ell_C-1)^2$; let $x\in[0,1)$ be the unique point with
$x(\ell_C-x)^2=4\varepsilon$.  The matrix
$\diag\bigl(x,\tfrac{\ell_C-x}{2},\tfrac{\ell_C-x}{2}\bigr)$
is positive semidefinite because $\ell_C>x$, has trace $\ell_C$ and determinant
$x(\ell_C-x)^2/4=\varepsilon$, and its least eigenvalue is $x<1$ since $3x<3<\ell_C$.
\end{proof}

\begin{proposition}[Lean]\label{prop:icosa}
At the icosahedral design of \S\ref{app:icosa} every triple $C$ satisfies
$\gamma(C)<0$, while the triple $\{0,2,4\}$ dominates strictly.
\end{proposition}

\begin{proof}
By \eqref{eq:icosapair} every leverage is $3$ and every distinct pairing has
square $\tfrac95$.  For a triple with pairings $p_{ab},p_{ac},p_{bc}$ and
product $\tau:=p_{ab}p_{ac}p_{bc}$,
\[
  \tr\Gram_C=9,
  \qquad
  \det\Gram_C=27+2\tau-3\bigl(p_{ab}^2+p_{ac}^2+p_{bc}^2\bigr)
             =27+2\tau-\tfrac{81}5=\tfrac{54}5+2\tau ,
\]
so by \eqref{eq:twomoment}
\[
  \gamma(C)=4\Bigl(\tfrac{54}5+2\tau\Bigr)-64=8\tau-\tfrac{104}5 .
\]
Now $\tau^2=(\tfrac95)^3=\tfrac{729}{125}$, while $(\tfrac{13}5)^2
=\tfrac{169}{25}=\tfrac{845}{125}$; hence $\lvert\tau\rvert<\tfrac{13}5$ and
$8\tau<\tfrac{104}5$, so $\gamma(C)<0$ at every one of the twenty triples.  The
domination of $\{0,2,4\}$, with $\lmin=3-3/\sqrt5>1$, is computed in
\S\ref{app:icosa}.
\end{proof}

The two-moment criterion is therefore silent at $(6,3)$, and by
Proposition~\ref{prop:momentsharp} so is every criterion in those two moments.
The criterion discards exactly the middle elementary symmetric function
of $\Gram_C$: at the icosahedron the moduli of the pairings are all equal, and
the verdict depends only on the sign of their product.  Any certificate for
$\Gtz(6,3)$ must consume $e_2(\Gram_C)$, or equivalently must be sensitive to
that sign.

\subsection{Structure inherited from smaller sizes}\label{app:structure}

\begin{remark}\label{rem:signature}
The signature of $W$ carries no information: by Lemma~\ref{lem:inertia}
the matrix $W=P-D$ has inertia $(k,m-k,0)$, and by Remark~\ref{rem:nogo}
there is a
Hermitian matrix of inertia $(1,1,0)$ with no positive semidefinite
$1\times1$ principal block.  So no statement of the form ``a Hermitian
matrix of inertia $(k,m-k,0)$ has a positive semidefinite $k\times k$
principal block'' is available; a proof must use more of the
projection than its inertia.
\end{remark}

By Corollary~\ref{cor:noparallel} the statement ``every extremal design
contains two parallel atoms'' is false at $(5,3)$: the diamond of
\S\ref{app:diamond} is extremal and its five atoms are pairwise non-parallel.
Every extremal design obtained by splitting does contain a parallel pair, so
the statement is true on the family one constructs first and false in general.
An induction branching on the presence of a parallel pair therefore owes a
separate treatment of the diamond, and of every extremal family not obtained by
splitting, a list not known to be finite.  This is part of
Problem~\ref{prob:structure}.

\subsection{Conditioning}\label{app:hygiene}

\begin{remark}\label{rem:hygiene}
The selection objective is ill-conditioned near the boundary of the weight
simplex.  As $t_c\to0$ the Parseval condition forces $\ell_c\sim1/t_c$, the
entries of $\SC$ grow like $1/\min_ct_c$, and the difference $\lmin(\SC)-1$ is
the difference of two quantities of that size.  A floating-point evaluation of
$\lmin$ on a formed Gram matrix therefore carries no information about the sign
of the margin; the stable quantity is $\sigma_{\min}(V_C)-1$, computed from the
isometry $V_C$ of \eqref{eq:chart} directly.  Every constant asserted in this
paper is exact: the witnesses above are rational, and the algebraic data of
Appendix~\ref{app:data} carries its defining relations.
\end{remark}
